\subsection{Conventions (ADR-0055)}
\label{sec:spn-0}

The binding conventions are recorded in
\texttt{docs/adr/0055-segmental-spine-kinematics-and-loads.md} (issue \#12736).

\paragraph{Per-Level Transform.} Each vertebra carries an ISB frame (Wu et al.\
2002): X anterior, Y cranial, Z right. Level $i$ maps the superior frame into the
inferior frame,
\begin{equation}
  T_i = T_{i,0}\, T\!\left(R_{ZXY}(\theta_i),\, d_i\right),
\end{equation}
with neutral offset $T_{i,0}$ and ISB joint-coordinate angles
$\theta_i = (\theta_{FE}, \theta_{LB}, \theta_{AR})$: extension, right lateral
bending and left axial rotation positive.

\paragraph{Composition.} For a chain to the named terminal body,
\begin{equation}
  T_{\mathrm{pelvis}\to\mathrm{terminal}} = T_{\mathrm{sacrum}}\, T_{L5S1}\, T_{L4L5} \cdots T_{\mathrm{terminal}},
\end{equation}
caudal factor first. Angles do not add: for three levels with
$(\theta_{FE},\theta_{LB},\theta_{AR})$ of $(-10,5,-15)$, $(-10,8,-20)$ and
$(-15,10,-25)$ degrees, the sum is $(-35, 23, -60)$ but the composed rotation
gives $(-25.36, 31.83, -57.57)$ degrees.

\paragraph{Loads.} Three observables (generalized effort, joint reaction with
static-optimization muscle forces, passive elements) are reported as the wrench
on the superior body by the inferior body, in the inferior ISB frame, about the
mid-disc centre. Compression is positive $+F_y$, AP shear $F_x$ (anterior),
ML shear $F_z$ (right).

\paragraph{Distribution.} Per-level angles from a coupling law solve
$\phi^\ast = \arg\min_\phi \lVert \log(R_{\mathrm{target}}^{\mathsf T} R(\phi))\rVert$
without altering the law; the geodesic residual is reported, and every output
carries its \texttt{distribution\_source}.

\paragraph{Limitations.} Per-level angles produced by a coupling law are model
assumptions, not measurements; golf literature load magnitudes are context
bands, not acceptance criteria.
